
\chapter{Introduzione e scopo del progetto}

La manipolazione fine, nota in letteratura come \textit{dexterous manipulation}, rappresenta una delle sfide più affascinanti e complesse nell'ambito della robotica moderna. 
Replicare la destrezza e la fluidità della mano umana richiede non solo una progettazione meccanica avanzata, ma anche algoritmi di cinematica, dinamica e controllo altamente sofisticati. 

Il presente progetto, denominato \textbf{RoboLiszt}, si pone l'obiettivo di progettare, modellare e simulare una mano robotica antropomorfa specificamente sviluppata per un compito di altissima precisione: l'esecuzione musicale al pianoforte.
La denominazione del progetto è stata coniata per simboleggiare la perfetta sinergia del trascendentale virtuosismo pianistico del celebre compositore ungherese Franz Liszt (1811-1886) con l'assoluta precisione deterministica propria dei moderni manipolatori robotici. 

\section{Stato dell'arte}

Lo sviluppo di RoboLiszt si inserisce in un filone di ricerca consolidato.
L'automazione applicata all'arte rappresenta da secoli uno dei banchi di prova più affascinanti per l'ingegneria.
L’interazione tra macchine e strumenti musicali non nasce con la moderna meccatronica: affonda le sue radici nell’illuminismo europeo, periodo in cui l’ambizione di replicare l’espressività umana portò alla nascita dei primi capolavori dell'orologeria.

Tra gli esempi storici più illustri spicca la \textit{Giocatrice di Timpano} (La Joueuse de Tympanon) di Pierre Kintzing e David Roentgen (1784), nonché i celebri automi di Pierre Jaquet-Droz, in particolare \textit{La Musicienne}, capace di suonare realmente un organo a canne premendo i tasti tramite un complesso sistema di camme meccaniche interne. 
In epoca contemporanea, progetti pionieristici come il robot pianista \textit{WABOT-2} (sviluppato dalla Waseda University negli anni '80) hanno segnato il passaggio definitivo dall'automa a movimento fisso al manipolatore governato da cinematiche digitali adattive.

\subsection{RoboPianist e il Reinforcement Learning}

Sviluppato da ricercatori di Google DeepMind, UC Berkeley e Stanford, \textit{RoboPianist}~\cite{robopianist2023} rappresenta l'attuale benchmark di riferimento per il controllo robotico ad alta dimensionalità. 
Il progetto utilizza il motore fisico MuJoCo per simulare due mani robotiche sospese sopra un pianoforte a 88 tasti. 

\begin{figure}[!ht]
    \centering
    \includegraphics[width=0.65\linewidth]{public/robopianist.png}
    \caption[In figura l'ambiente di simulazione per l'esecuzione pianistica di RoboPianist utilizzando il motore fisico open-source MuJoCo.]{In figura l'ambiente di simulazione per l'esecuzione pianistica di RoboPianist utilizzando il motore fisico open-source MuJoCo. 
    Esso è costituito da una tastiera digitale a grandezza naturale da 88 tasti e da due mani robotiche Shadow Dexterous Hands, ciascuna dotata di 24 gradi di libertà.
    Fonte:~\cite{robopianist2023}.}
    \label{fig:robopianist}
\end{figure}

A differenza di RoboLiszt, che impiega un approccio analitico basato sulla cinematica inversa e il controllo dinamico, RoboPianist sfrutta il \textit{Deep Reinforcement Learning} per insegnare a un agente autonomo l'esecuzione di partiture complesse in formato MIDI, nello specifico l'algoritmo DroQ.

Dal punto di vista  meccanico, il principale riferimento accademico è la \textit{Shadow Dexterous Hand} (modello E3M5). 
Si tratta di una delle mani robotiche più all'avanguardia attualmente in commercio: caratterizzata da dimensioni umane, replica la complessa biomeccanica dei tendini grazie a un sofisticato sistema a 24 gradi di libertà per ogni singola mano.
Pur essendo un hardware fisico, i suoi modelli cinematici open-source costituiscono lo standard per la simulazione e forniscono i parametri anatomici realistici (quali lunghezze delle falangi, offset palmari e limiti articolari) su cui si basa l'albero cinematico sviluppato in questo progetto.

Il brano musicale viene codificato partendo dallo standard MIDI e convertito in una traiettoria temporale (\textit{piano roll}), che funge da stato obiettivo per l'agente. 
Ad ogni istante di tempo, la rete neurale riceve in ingresso lo stato propriocettivo della mano e lo stato dei tasti, restituendo come output gli angoli di giunto target per gli attuatori proporzionali.

Una delle principali sfide evidenziate dallo studio è la difficoltà di esplorazione dello spazio delle azioni a causa dell'altissima dimensionalità del sistema. 
Per ovviare a questo problema ed evitare segnali di \textit{reward} troppo sparsi, i ricercatori hanno integrato euristiche umane all'interno della funzione di ricompensa, utilizzando un dataset di diteggiature pianistiche (\textit{Repertoire-150}) per guidare l'apprendimento dell'agente. 
Infine, l'accuratezza dell'esecuzione viene misurata tramite metriche statistiche standard, quali \textit{Precision}, \textit{Recall} e \textit{F1 score}, per valutare l'esatta sincronizzazione e pressione dei tasti richiesti dalla partitura originale.

\subsection{Apprendimento autonomo: il caso della USC Viterbi}

Un ulteriore paradigma innovativo nel campo della \textit{dexterous manipulation} musicale è stato recentemente illustrato dai ricercatori della Viterbi School of Engineering presso la University of Southern California (USC)~\cite{usc_pianohand_2026}. 
Lo studio presenta una mano robotica in grado di apprendere autonomamente l'esecuzione di brani al pianoforte senza fare affidamento su rigide traiettorie pre-programmate. 

A differenza dei robot convenzionali basati sull'inversione dello Jacobiano e su modelli dinamici espliciti (come l'architettura proposta in \textit{RoboLiszt}), l'approccio della USC si fonda su algoritmi di apprendimento automatico. 
Il sistema è progettato per auto-istruirsi (\textit{self-taught}) interagendo direttamente con lo strumento: valutando i propri errori acustici e di posizionamento, la mano affina iterativamente la coordinazione e la fluidità dei movimenti. 

La mano utilizza quattro dita azionate da tendini e controllate da piccoli motori elettrici, progettati per imitare la meccanica della mano umana. 
Inoltre, un sistema di reti neurali analizzano il suono della melodia e lo convertono nei comandi motori necessari per riprodurla.

Il robot pianista rappresenta una dimostrazione di fattibilità per quella che i ricercatori chiamano “robotica percettiva”: un modello in cui il sistema percepisce l'ambiente circostante, sperimenta movimenti e si autocorregge senza ricorrere a dati di addestramento esaustivi. 
Secondo gli studiosi, questo stesso approccio potrebbe in futuro offrire assistenza alle persone in modi molto più personali e intuitivi rispetto a quanto facciano gli attuali robot progettati per compiti specifici.

\begin{figure}[!ht]
    \centering
    \includegraphics[width=0.5\linewidth]{public/robohand.png}
    \caption[Melodia tratta da \textit{Perception in action}: un sistema robotico in grado di imparare autonomamente a suonare delle melodie musicali.]{Melodia tratta da \textit{Perception in action}: un sistema robotico in grado di imparare autonomamente a suonare delle melodie musicali.
    Fonte:~\cite{usc_pianohand_2026}.}
    \label{fig:robohand}
\end{figure}

\subsection{Apprendimento data-driven per task atletici dinamici}

A dimostrazione di come le moderne architetture di controllo si stiano espandendo verso ogni dominio fisico, un recente e rilevante contributo è presentato in uno studio pubblicato su arXiv~\cite{athletic_tasks_2025}. 
Il progetto esplora lo sviluppo di sistemi robotici in grado di eseguire task atletici altamente dinamici, ciascuno caratterizzato da sfide fisiche e requisiti di agilità unici. 

\begin{figure}[!ht]
    \centering
    \includegraphics[width=0.75\linewidth]{public/rp1m}
    \caption[Agente basato sull'apprendimento per rinforzo (RL) che elimina la necessità di diteggiature annotate da esseri umani, formulando il posizionamento delle dita come un problema di Trasporto Ottimale (OT).]{L'articolo~\cite{athletic_tasks_2025} propone un agente basato sull'apprendimento per rinforzo (RL) che elimina la necessità di diteggiature annotate da esseri umani, formulando il posizionamento delle dita come un problema di Trasporto Ottimale (OT); ciò consente all'agente di eseguire brani musicali senza alcun aiuto umano, incluso il complesso brano “Il volo del calabrone”. 
    Viene inoltre realizzato un dataset su larga scala, denominato RP1M++, costituito da oltre un milione di traiettorie di esecuzione pianistica esperta, ottenute addestrando più di 2$\,$000 agenti RL.}
    \label{fig:rp1m}
\end{figure}

Dal punto di vista dell'infrastruttura software, framework di questa portata fanno tipicamente affidamento su un'architettura ibrida, utilizzando \textbf{Python} per l'elaborazione dei dati e il \textit{training} dei modelli di apprendimento (sfruttando librerie come PyTorch o JAX), accoppiato al \textbf{C++} per garantire le frequenze di calcolo necessarie al deployment del controllo in \textit{real-time} sull'hardware.

L'aspetto più ragguardevole di questo progetto risiede nella capacità dell'algoritmo di “consumare” e generalizzare enormi moli di dati per gestire vincoli dinamici severi. 
Questo approccio offre un interessante spunto di riflessione se confrontato con \textit{RoboLiszt}: mentre la mano pianista opera nel dominio della manipolazione fine (richiedendo estrema precisione e destrezza su microscala), il robot atletico affronta il problema duale, dovendo governare forze impulsive, inerzie elevate e stabilità globale su macroscala. 
Entrambi i domini, tuttavia, evidenziano la necessità di modelli dinamici accurati e algoritmi di controllo robusti.

\subsection{Un approccio bio-ispirato}

Un approccio profondamente diverso dalla modellazione a giunti rigidi è rappresentato dal sistema robotico descritto in un recente studio pubblicato sul \textit{Journal of the Royal Society Interface}~\cite{perception_action_2025}. 
La ricerca illustra il controllo dinamico di una mano robotica, definita \textit{Musician Hand}, dotata di quattro dita attuate a tendine e guidate da altrettanti motori in corrente continua (DC).

Come dettagliato nella Sezione~2.2 dello studio, la meccanica fisica del manipolatore emula l'anatomia umana. 
Ogni dito è modellato rigorosamente come un sistema dinamico del secondo ordine, la cui equazione del moto integra l'inerzia articolare ($I$), le forze centrifughe e di Coriolis ($c$), il coefficiente di attrito al giunto ($b$) e la rigidezza strutturale.
In formule:
\begin{equation}
    \ddot{q}_i =
    -I(q_i)^{-1} c(q_i, \dot{q}_i) + b \dot{q}_i + I(q_i)^{-1} T
    \label{eq:musician_hand_meccanics}
\end{equation}
La coppia motrice non è applicata direttamente sull'asse di rotazione, ma è generata dalle forze “muscolo-tendinee” (i cavi trazionati dai motori) moltiplicate per il rispettivo braccio di momento meccanico. 

\begin{figure}[!ht]
    \centering
    \includegraphics[width=0.95\linewidth]{public/musician_hand}
    \caption[Profili di attivazione motoria della \textit{Musician Hand} e sequenze di attivazione motoria.]{(a) Profili di attivazione motoria della \textit{Musician Hand}: per la riproduzione delle melodie sono stati definiti quattro profili di attivazione pre-programmati. 
    La funzione di controllo impiega un'approssimazione esponenziale della dinamica di attivazione-contrazione muscolare [50--52]. 
    Per emulare il fisiologico tono muscolare e garantire la prontezza operativa, ciascun profilo è inizializzato al 5\% dell'attivazione massima, raggiungendo successivamente quattro distinti livelli di plateau dell'intensità ($K_{1-4}$) regolati da specifiche costanti di tempo ($\tau_{1-4}$). 
    Tale approccio genera quattro differenti meccaniche di interazione con la tastiera e da vita a variazioni precise nell'attacco, nel volume e nel decadimento acustico della nota. 
    (b) Sequenze di attivazione motoria: esempio dei segnali di controllo inviati agli attuatori per la trazione dei tendini e la conseguente pressione dei tasti.
    Fonte:~\cite{perception_action_2025}.}
    \label{fig:musician_hand}
\end{figure}

L'aspetto più innovativo del progetto risiede nel paradigma di attuazione: il sistema mappa direttamente lo spettrogramma dell'audio musicale in serie temporali di attuazione tramite una Rete Neurale Artificiale (costituita da CNN e MLP). 
I motori vengono azionati seguendo un'approssimazione esponenziale della dinamica di attivazione-contrazione muscolare umana, impiegando quattro differenti plateau di intensità e costanti di tempo per replicare realisticamente l'attacco, il volume e il decadimento acustico della nota.

\subsection{Pianobot: la pianificazione della traiettoria}

Un ulteriore e pragmatico contributo allo stato dell'arte è offerto dai framework \textit{open-source} orientati all'integrazione di sistema, tra cui spicca il progetto \textit{Pianobot}~\cite{pianobot_lin}. 
A differenza delle complesse architetture bio-ispirate o basate su reti neurali, questo repository mette in luce l'implementazione algoritmica pura, poiché fornisce una pipeline software strutturata per la traduzione diretta di file MIDI in traiettorie cinematiche eseguibili.

Sebbene non si concentri sull'ottimizzazione anatomica e sulla riduzione mirata dei gradi di libertà che caratterizza \textit{RoboLiszt}, \textit{Pianobot} rappresenta un eccellente riferimento architetturale per la pianificazione computazionale del moto (\textit{trajectory planning}) e per la rigorosa temporizzazione (\textit{timing}) delle sequenze musicali.
Tra le funzionalità implementate vi troviamo l'utilizzo di algoritmi di Computer Vision fondamentali per l'individuazione del pianoforte, come ad esempio l'algoritmo SIFT (\textit{Scale-Invariant Feature Transform}), indispensabile per la localizzazione dello strumento.

Un aspetto particolarmente affine alle problematiche affrontate in \textit{RoboLiszt} è la successiva manipolazione algoritmica delle note: attraverso specifiche tecniche di segmentazione e \textit{clustering}, il framework è in grado di calcolare il percorso spaziale più breve per lo scorrimento del braccio robotico e di determinare la sequenza ottima per il movimento delle singole dita.

\section{Motivazione della scelta}

Il presente progetto nasce dalla volontà di coniugare due discipline all'apparenza distanti, ma intrinsecamente legate da profonde affinità: l'informatica, declinata nei campi della robotica e del controllo dei sistemi, e la musica. 
Sebbene operino in domini differenti, entrambe le materie condividono un fondamento basato sul rigore matematico, sulla logica strutturale, sulla ritmica e sulla precisione algoritmica.

La realizzazione di un sistema cibernetico in grado di interfacciarsi fisicamente con uno strumento musicale rappresenta la sintesi perfetta di tali elementi. 
In questa prospettiva, la partitura musicale smette di essere un semplice costrutto artistico e diventa una rigorosa pianificazione di traiettorie nel tempo e nello spazio operativo, governata da precise leggi di cinematica e dinamica.

\begin{figure}[!ht]
    \centering
    \includegraphics[width=0.75\linewidth]{public/robotic_hand_plays_piano.jpeg}
    \caption{Rappresentazione della mano robotica al pianoforte.}
\end{figure}

\subsection{L'esecuzione pianistica come benchmark per la robotica}

Il pianoforte è uno dei compiti più esigenti per una mano, umana o artificiale. 
L'esecuzione richiede un posizionamento spaziale preciso (i tasti bianchi hanno una larghezza di circa 23.5~mm), una coordinazione simultanea di più dita (accordi), movimenti rapidi del polso e dell'avambraccio per coprire l'estensione della tastiera, e un controllo del \emph{tocco}, cioè della forza e della velocità con cui il tasto viene colpito, da cui dipende l'intensità del suono. 
Il pianoforte è un banco di prova naturale per lo studio della cinematica, della dinamica e del controllo di mani robotiche antropomorfe, specialmente per diverse ragioni:
\begin{itemize}
    \item \textbf{Spazio di lavoro vincolato e discreto:} I tasti di un pianoforte sono geometricamente noti, complanari e disposti a passo costante. Ciò consente di formulare target cartesiani deterministici ($x_k$, $y_k$, $z_k$), ideali per la verifica analitica della cinematica inversa.
    \item \textbf{Manipolazione a contatto intermittente:} a differenza dei robot di saldatura o verniciatura a traiettoria continua, il tocco del tasto alterna fasi di moto libero (spostamento e posizionamento) a fasi di contatto verticale rapido con precisione millimetrica.
    \item \textbf{Cinematica ramificata a catene aperte:} permette di studiare l'estensione delle catene seriali verso strutture ad albero, senza incorrere nelle complicazioni dinamiche e vincolari delle catene chiuse.
\end{itemize}
L'insieme di queste specifiche sfide cinematiche costituirà il fulcro dell'analisi condotta nei capitoli successivi del presente lavoro, al fine di validare rigorosamente le scelte architetturali e di controllo adottate nel sistema.

\subsection{Obiettivi del progetto}

Il presente lavoro si propone di:
\begin{itemize}
    \item definire un modello cinematico a catena aperta ramificata della mano (polso a 2 DoF, quattro dita e pollice, per un totale di 20 giunti rotoidali) e analizzarne spazio di lavoro, singolarità e ridondanza;
    \item sviluppare un algoritmo di cinematica inversa robusto alle singolarità e ai limiti articolari, capace di coordinare più dita attorno a un polso condiviso;
    \item caratterizzare la capacità della struttura di trasmettere velocità e forze al polpastrello tramite le ellissoidi di manipolabilità, in forma cinematica e dinamica;
    \item formulare la dinamica della struttura ramificata con la formulazione ricorsiva di Newton--Eulero e utilizzarla per calcolare le coppie necessarie lungo i task;
    \item progettare un controllo ibrido posizione/forza per l'interazione con la tastiera;
    \item realizzare un'architettura software che colleghi modello, pianificazione e simulazione, fino alla simulazione fisica in ambiente Gazebo.
\end{itemize}
